%------------------------------------------------------------------------------%
\documentclass[a4paper,12pt,fleqn]{article}

\usepackage{calculo_varias_variaveis-1}

\ShowAnswers

%------------------------------------------------------------------------------%
\begin{document}

\makeHeader{CFVVI}{Prova 3 -- Turma 8}{20/05/2026}

\bigskip
\textbf{Atenção}: a aplicação das regras do produto ou do quociente em casos
nos quais um dos termos é constante \textbf{será considerada incorreta},
por evidenciar compreensão inadequada das regras de derivação.

% 01
%------------------------------------------------------------------------------%
\exercise{50\%}
Considere a função
\(
  g(x, y, z) = \sqrt{x+y^2}\,\cos(yz)+\frac{x}{y+z}+z^3\ln(x+1)
\)
\begin{enumerate}[label=\alph*),nosep]
  \item Calcule o gradiente de $f$

  \item Calcule a derivada direcional de $f$ no ponto
  \(
    Q = (0, 1, 1)
  \)
  na direção
  \(
    w = (1, 2, -2)
  \)

  \item Determine a direção em que $f$ cresce mais rapidamente no ponto $P$
\end{enumerate}

\begin{answer}
  \textbf{(a)} O gradiente de
  \(
    g(x, y, z)=\sqrt{x+y^2}\,\cos(yz)+\frac{x}{y+z}+z^3\ln(x+1)
  \)
  é
  \(
    \nabla f(x, y, z)
    = \Vetortri{\D{g}{x}}{\D{g}{y}}{\D{g}{z}}
  \)
  \begin{align*}
    \D{g}{x}
    & = \D{}{x} \left( e^{xy}\ln\left(x+z^2\right) \right)
      + \D{}{x} \left( \frac{y^2}{z} \right)
      - \D{}{x} \left( \sin(xz) \right)
    \\[2mm]
    & = ye^{xy}\ln\left(x+z^2\right)
      + e^{xy}\frac{1}{x+z^2}
      - z\cos(xz)
  \end{align*}
  \begin{align*}
    \D{f}{y}
    & = \D{}{y} \left( e^{xy}\ln\left(x+z^2\right) \right)
      + \D{}{y} \left( \frac{y^2}{z} \right)
      - \D{}{y} \left( \sin(xz) \right)
    \\[2mm]
    & = xe^{xy}\ln\left(x+z^2\right)
      + \frac{2y}{z}
  \end{align*}
  \begin{align*}
    \D{f}{z}
    & = \D{}{z} \left( e^{xy}\ln\left(x+z^2\right) \right)
      + \D{}{z} \left( \frac{y^2}{z} \right)
      - \D{}{z} \left( \sin(xz) \right)
    \\[2mm]
    & = e^{xy}\frac{2z}{x+z^2}
      - \frac{y^2}{z^2}
      - x\cos(xz)
  \end{align*}
  \[
    \nabla f(x,y,z)
    = \Vetortri{
    ye^{xy}\ln(x+z^2)
    +
    \frac{e^{xy}}{x+z^2}
    -
    z\cos(xz)}{
    xe^{xy}\ln(x+z^2)
    +
    \frac{2y}{z}}{
    \frac{2ze^{xy}}{x+z^2}
    -
    \frac{y^2}{z^2}
    -
    x\cos(xz)}
  \]

  \textbf{(b)} No ponto \(P=(1, 1, 1)\),
  \[
    \nabla f(1, 1, 1)
    = \Vetortri{e\ln2+\frac{e}{2}-\cos1}{e\ln2+2}{e-1-\cos1}
  \]
  O vetor unitário na direção de
  \(
    v = (2, -1, 2)
  \)
  é
  \[
    u
    = \frac{v}{\norm{v}}
    = \frac{1}{3} \vetortri{2}{-1}{2}
    = \Vetortri{\frac{2}{3}}{-\frac{1}{3}}{\frac{2}{3}}
  \]
  A derivada direcional é
  \begin{align*}
    D_{u} f(1, 1, 1)
    & = \nabla f(1, 1, 1) \cdot u                                        \\[2mm]
    & = \Vetortri{e\ln2+\frac{e}{2}-\cos1}{e\ln2+2}{e-1-\cos1} \cdot
        \Vetortri{\frac{2}{3}}{-\frac{1}{3}}{\frac{2}{3}}                \\[2mm]
    & = \frac{2}{3} \left( e\ln 2 + \frac{e}{2} - \cos(1) \right)
      - \frac{1}{3} \left( e\ln 2 + 2 \right)
      + \frac23 \left( e - 1 - \cos(1) \right)                           \\[2mm]
    & = \frac{1}{3} \left( e\ln 2 + 3e - 4\cos(1) - 4 \right)
  \end{align*}

  \textbf{(c)} A função cresce mais rapidamente na direção do gradiente
  \[
    \nabla f(1, 1, 1)
    =\Vetortri{e\ln2+\frac{e}{2}-\cos1}{e\ln2+2}{e-1-\cos1}
  \]
\end{answer}

% 02
%------------------------------------------------------------------------------%
\exercise{50\%}
Considerando a função
\(
  f(x, y) = x^2 e^y + \ln(x+y)
\)
\begin{enumerate}[label=\alph*),nosep]
  \item Determine a equação do plano tangente ao gráfico de $f$
        no ponto correspondente a \((1, 0)\)
  \item Verifique se o ponto \((-1, 2, 1)\) pertence ao plano tangente
\end{enumerate}

\begin{answer}

  \textbf{(a)}

  A equação do plano tangente ao gráfico de \(z=f(x,y)\) no ponto
  \[
    (x_0,y_0,z_0)
  \]
  é dada por
  \[
    z-z_0
    = f_x(x_0,y_0)(x-x_0)
    + f_y(x_0,y_0)(y-y_0).
  \]
  
  Primeiramente, calculamos o valor da função:
  
  \[
  f(1,0)
  =
  1^2e^0+\ln(1+0)
  =
  1.
  \]
  
  Logo,
  
  \[
  z_0=1.
  \]
  
  Calculamos agora as derivadas parciais.
  
  Para \(x\),
  
  \[
  f_x(x,y)
  =
  2xe^y+\frac{1}{x+y}.
  \]
  
  Para \(y\),
  
  \[
  f_y(x,y)
  =
  x^2e^y+\frac{1}{x+y}.
  \]
  
  Avaliando em \((1,0)\),
  
  \[
  f_x(1,0)
  =
  2(1)e^0+\frac{1}{1}
  =
  3,
  \]
  
  \[
  f_y(1,0)
  =
  1^2e^0+\frac{1}{1}
  =
  2.
  \]
  
  Substituindo na fórmula do plano tangente,
  
  \[
  z-1
  =
  3(x-1)
  +
  2(y-0).
  \]
  
  Simplificando,
  
  \[
  z
  =
  3x+2y-2.
  \]
  
    \bigskip
    \textbf{(b)}
    O ponto não pertence ao plano, pois
    \begin{align*}
      z
      & = \left[3x+2y-2\right]\evalat{(-1, 2)}{} \\
      & = -3 +4 -2 \\
      & = -1 \\
      & \neq 1
    \end{align*}
  
  
  
  \hrulefill
  
  \hrulefill
  
    \textbf{(a)}
    A equação do plano tangente ao gráfico de \(z = g(x, y)\) no ponto
    \(
      (a, b, c)
    \)
    é dada por
    \[
      z - c = g_x(a, b)(x - a) + g_y(a, b)(y - b)
    \]
    \[
      c
      = g(0, \pi)
      = \sqrt{1} + \sin(-\pi)
      = 1
    \]
    Derivadas parciais
    \[
      \D{g}{x}
      = \frac{y}{2\sqrt{xy+1}} + \cos(x-y)
    \]
    \[
      \D{g}{y}
      = \frac{x}{2\sqrt{xy+1}} - \cos(x-y)
    \]
    Avaliando em \((0, \pi)\)
    \[
      \D{g}{x}(0, \pi)
      = \frac{\pi}{2\sqrt{0+1}} + \cos(-\pi)
      = \frac{\pi}{2} - 1
    \]
    \[
      \D{g}{y}(0, \pi)
      = \frac{0}{2\sqrt{0+1}} - \cos(0-\pi)
      = 1
    \]
    Substituindo na fórmula
    \[
      z - 1
      = \left( \frac{\pi}{2} - 1 \right)(x-0)
      + 1 (y-\pi)
    \]
    \[
      z - 1
      = \left( \frac{\pi}{2} - 1 \right)x
      + y -\pi
    \]
    \[
      z
      = \left( \frac{\pi}{2} - 1 \right)x
      + y -\pi +1
    \]
  
    \bigskip
    \textbf{(b)}
    O ponto não pertence ao plano, pois
    \begin{align*}
      z
      & = \left[\left( \frac{\pi}{2} - 1 \right)x + y -\pi +1\right]\evalat{(2, 0)}{} \\
      & = \left( \frac{\pi}{2} - 1 \right)2 + 0 -\pi +1 \\
      & = \pi - 2 -\pi +1 \\
      & = -1 \\
      & \neq 1
    \end{align*}
\end{answer}

%------------------------------------------------------------------------------%
\end{document}
%------------------------------------------------------------------------------%
